% =====================================================================
%  PREAMBOLO per la classe unicam-diss (involucro di tufte-book v3.5.2)
%
%  Da inserire nel main del template DOPO \documentclass, al posto del
%  preambolo di main.tex (che era scritto per la classe "report").
%
%  Regola: la classe carica gia' geometry, setspace, fancyhdr, hyperref,
%  xcolor, inputenc, fontenc, titlesec, titletoc, ragged2e, paralist,
%  multicol, tocloft e i font (ccfonts + XCharter + beramono).
%  NON ricaricarli: tufte-book li configura e un secondo \usepackage con
%  opzioni diverse causa "Option clash for package ...".
%  Elenco verificato sul log di compilazione del template (21/08/2026).
% =====================================================================

% ---------- Lingua ----------
% inputenc/fontenc sono gia' caricati dalla classe: qui serve solo babel.
\usepackage[italian]{babel}
\usepackage{csquotes}

% ---------- Matematica ----------
\usepackage{amsmath,amssymb}

% ---------- Tabelle ----------
\usepackage{booktabs}
\usepackage{tabularx}
\usepackage{array}
\usepackage{multirow}
% NOTA: longtable e' compatibile con tufte solo fuori dal blocco di testo
% stretto; se serve, usarlo dentro l'ambiente fullwidth della classe.

% ---------- Elenchi ----------
% ATTENZIONE: la classe carica gia' "paralist", che ridefinisce gli elenchi.
% Il documento usa 46 elenchi con opzioni enumitem (leftmargin, itemsep,
% label). Se enumitem entra in conflitto, le alternative in ordine di costo:
%   1. caricare enumitem PRIMA di questo blocco e verificare;
%   2. rimuovere le opzioni e usare gli elenchi nativi;
%   3. definire una volta gli elenchi con \setlist e togliere le opzioni
%      dai singoli \begin{itemize}/\begin{enumerate}.
\usepackage{enumitem}

% ---------- Simboli ----------
\usepackage{gensymb}   % \degree, \celsius
\usepackage{eurosym}   % \euro

% ---------- Codice ----------
% tufte-book non gestisce i listati: listings si carica normalmente.
% I listati vanno tenuti stretti (la colonna di testo di tufte e' circa 2/3
% della pagina) oppure messi in fullwidth.
\usepackage{listings}
\lstset{
  basicstyle=\ttfamily\scriptsize,
  breaklines=true,
  showstringspaces=false,
  frame=single,
  framesep=4pt,
  captionpos=b
}

% ---------- Riferimenti ----------
% hyperref e' gia' caricato dalla classe: cleveref va DOPO, e basta lui.
\usepackage{cleveref}

% ---------- Bibliografia ----------
% Il template usa gia' biblatex con backend biber (confermato dai file
% output.bcf/output.blg): NON ricaricare biblatex, cambiare solo la risorsa.
% \addbibresource{bib/tesi.bib}

% ---------- Comandi di comodo del documento ----------
\newcommand{\ld}{LD2410B}
\newcommand{\ldb}{HLK-LD2410B}
\newcommand{\ldc}{HLK-LD2420}
\newcommand{\pir}{HC-SR501}
% Note redazionali di bozza: per la stampa finale svuotare la macro.
\newcommand{\todo}[1]{\textcolor{red}{\textbf{[TODO: #1]}}}
% \renewcommand{\todo}[1]{}

% =====================================================================
%  Da NON ricaricare (gia' nella classe): geometry, setspace, fancyhdr,
%  hyperref, xcolor, inputenc, fontenc, titlesec, titletoc, ragged2e,
%  paralist, multicol, tocloft, textcomp, placeins, changepage.
%  Da NON caricare affatto con tufte: caption, subcaption, emptypage.
%  tufte-book ha un proprio sistema di didascalie e di note a margine;
%  "caption" lo rompe.
% =====================================================================
